%
% Para resaltar secciones Revisadas
% \begin{revisado}
% ///
% \end{revisado}
%
% Para resaltar secciones a Corregir
% \begin{corregir}
% ///
% \end{corregir}
%
\usepackage{tcolorbox}
\tcbuselibrary{breakable}
\definecolor{celeste}{rgb}{0.7,0.9,1}
\tcbset{
    revisionbox/.style={
        boxrule=0pt,
        sharp corners,
        breakable,
        before upper={\setlength{\parindent}{15pt}}
    }
}
\newtcolorbox{revisado}{revisionbox, colback=celeste!50}
\newtcolorbox{corregir}{revisionbox, colback=yellow!50}
\newtcbox{\inlineentitybox}[1][]{%
    on line,
    boxrule=0.8pt,
    arc=2pt,
    boxsep=1pt,
    left=3pt,
    right=3pt,
    top=1pt,
    bottom=1pt,
    #1
}
\newcommand{\entitytag}[3]{%
    \inlineentitybox[colframe=#1!75!black,colback=#1!10]{\strut #3}%
    \textsuperscript{\textcolor{#1!75!black}{\textsf{#2}}}%
}
\newcommand{\entityairportcode}[1]{\entitytag{blue}{AIRPORT\_CODE}{#1}}
\newcommand{\entityflightnumber}[1]{\entitytag{cyan}{FLIGHT\_NUMBER}{#1}}
\newcommand{\entitypnrcode}[1]{\entitytag{violet}{PNR\_CODE}{#1}}
\newcommand{\entityperson}[1]{\entitytag{red}{PERSON}{#1}}
\newcommand{\entitydate}[1]{\entitytag{orange}{DATE}{#1}}
\newcommand{\entitytime}[1]{\entitytag{green}{TIME}{#1}}
\newcommand{\entityflightsegmentmarker}[1]{\entitytag{yellow}{TIME}{#1}}
\newcommand{\entityairportname}[1]{\entitytag{teal}{AIRPORT\_NAME}{#1}}
\newcommand{\entityairlinename}[1]{\entitytag{magenta}{AIRLINE\_NAME}{#1}}
\newcommand{\entitystateorprovince}[1]{\entitytag{black}{STATE\_OR\_PROVINCE}{#1}}
\newcommand{\entityorganization}[1]{\entitytag{black}{ORGANIZATION}{#1}}
\newcommand{\entitynumber}[1]{\entitytag{black}{NUMBER}{#1}}
\newcommand{\entityduration}[1]{\entitytag{black}{DURATION}{#1}}

\setcounter{tocdepth}{5}
\setcounter{secnumdepth}{6}

\usepackage{geometry}                % See geometry.pdf to learn the layout options. There are lots.
\geometry{a4paper}
\usepackage[T1]{fontenc}
\usepackage{lmodern}
\usepackage{graphicx}
\usepackage{tikz}
\usepackage{colortbl}
\usepackage{fvextra}
\usepackage{listings}
\usepackage{csquotes}
\usepackage{microtype}
\usetikzlibrary{arrows.meta,fit,positioning,shapes.multipart}

\setlength{\parskip}{0.25\baselineskip} % Espacio entre párrafos
\setlength{\parindent}{15pt} % Sangría en la primera línea de cada párrafo

\usepackage[spanish,es-noshorthands]{babel}
\addto\captionsspanish{%
    \renewcommand{\figureautorefname}{Figura}%
    \renewcommand{\tableautorefname}{Tabla}%
    \renewcommand{\sectionautorefname}{Sección}%
    \renewcommand{\subsectionautorefname}{Subsección}%
    \renewcommand{\subsubsectionautorefname}{Subsubsección}%
}
\usepackage[
    backend=biber,
    style=apa,
    sorting=nyt,
    language=spanish,
    maxcitenames=99,
    mincitenames=99
]{biblatex}
\DeclareLanguageMapping{spanish}{spanish-apa}
\DefineBibliographyStrings{spanish}{
    andothers = {et\addabbrvspace al},
    urlseen = {Accedido}
}
% Para sitios web cargados como @misc, preferimos citar el título sin comillas.
\DeclareFieldFormat[misc]{citetitle}{#1}
\newcommand{\legacyurlmonthabbr}[1]{%
    \ifcase#1\or Jan\or Feb\or Mar\or Apr\or May\or Jun\or Jul\or Aug\or Sep\or Oct\or Nov\or Dec\fi}
\newcommand{\legacyurldate}{%
    \mkbibparens{Accedido:\space
    \iffieldundef{urlday}{}{\thefield{urlday}}%
    \iffieldundef{urlmonth}{}{-\legacyurlmonthabbr{\thefield{urlmonth}}}%
    \iffieldundef{urlyear}{}{-\thefield{urlyear}}}}
\renewbibmacro*{url+urldate}{%
    \ifthenelse{\iffieldundef{url}\OR\NOT\iffieldundef{doi}}
    {}
    {\printtext{Recuperado de\space\printfield{url}}%
    \iffieldundef{urlyear}
    {}
    {\space\legacyurldate}}}
\DeclareNameWrapperFormat{citehyperref}{\bibhyperref{#1}}
\DeclareNameWrapperAlias{labelname}{citehyperref}
\DeclareDelimFormat[textcite]{nameyeardelim}{\addspace}
\DeclareDelimFormat[cite]{nameyeardelim}{\addspace}
\makeatletter
\DeclareCiteCommand{\cite}[\cbx@textcite@init\cbx@textcite]
{\gdef\cbx@savedkeys{}%
\citetrackerfalse%
\pagetrackerfalse%
\DeferNextCitekeyHook%
\usebibmacro{cite:init}}
{\ifthenelse{\iffirstcitekey\AND\value{multicitetotal}>0}
{\protected@xappto\cbx@savedcites{()(\thefield{multipostnote})}%
\global\clearfield{multipostnote}}
{}%
\xappto\cbx@savedkeys{\thefield{entrykey},}%
\iffieldequals{namehash}{\cbx@lasthash}
{}
{\stepcounter{textcitetotal}%
\savefield{namehash}{\cbx@lasthash}}}
{}
{\protected@xappto\cbx@savedcites{%
    [\thefield{prenote}][\thefield{postnote}]{\cbx@savedkeys}}}
\makeatother
\addbibresource{references.bib}

\usepackage{titlesec}
\titlelabel{\thetitle\quad}  % Esto elimina el punto

\usepackage{xurl}
\usepackage[hidelinks,hypertexnames=false]{hyperref} % Para autoref
\hypersetup{bookmarksdepth=5}

\lstdefinestyle{basecodestyle}{
    basicstyle=\ttfamily\small,
    breaklines=true,
    breakatwhitespace=false,
    breakautoindent=true,
    breakindent=0pt,
    columns=fullflexible,
    keepspaces=true,
    showstringspaces=false,
    captionpos=b,
    frame=single,
    framerule=0.4pt,
    rulecolor=\color{black},
    upquote=true,
    literate=
        {á}{{\'a}}1 {é}{{\'e}}1 {í}{{\'\i}}1 {ó}{{\'o}}1 {ú}{{\'u}}1
        {Á}{{\'A}}1 {É}{{\'E}}1 {Í}{{\'I}}1 {Ó}{{\'O}}1 {Ú}{{\'U}}1
        {ñ}{{\~n}}1 {Ñ}{{\~N}}1
        {ü}{{\"u}}1 {Ü}{{\"U}}1
        {ö}{{\"o}}1 {Ö}{{\"O}}1
        {–}{{--}}1
        {“}{{``}}1 {”}{{''}}1 {’}{{'}}1
}

\lstdefinestyle{javacodestyle}{
    style=basecodestyle,
    language=Java,
    keywordstyle=\color{blue!70!black}\bfseries,
    stringstyle=\color{green!40!black},
    commentstyle=\color{gray!70!black}\itshape
}

\lstdefinelanguage{jsonlike}{
    morestring=[b]",
    alsoletter={\$},
    morekeywords={ner,type,value,ruleType,pattern,action,result,tokens,Annotate},
    sensitive=true
}

\lstdefinestyle{jsonlikecodestyle}{
    style=basecodestyle,
    language=jsonlike,
    keywordstyle=\color{blue!70!black}\bfseries,
    stringstyle=\color{green!40!black}
}

\lstnewenvironment{javacode}{\lstset{style=javacodestyle}}{}
\lstnewenvironment{jsonlikecode}{\lstset{style=jsonlikecodestyle}}{}
\renewcommand{\lstlistingname}{Fragmento}
\providecommand*{\lstlistingautorefname}{Fragmento}

\DefineVerbatimEnvironment{codeblock}{Verbatim}{
    fontsize=\small,
    breaklines=true,
    frame=single,
    framesep=2mm
}

\titleformat{\section}
{\normalfont\Large\bfseries}
{\thesection}{1em}{}

\titleformat{\subsection}
{\normalfont\large\bfseries}
{\thesubsection}{1em}{}

\titleformat{\subsubsection}
{\normalfont\normalsize\bfseries}
{\thesubsubsection}{1em}{}

\titleformat{\paragraph}[block]
{\normalfont\normalsize\bfseries}
{\theparagraph}{1em}{}

\titleformat{\subparagraph}[block]
{\normalfont\normalsize\bfseries}
{\thesubparagraph}{1em}{}

\titleclass{\subsubparagraph}{straight}[\subparagraph]
\newcounter{subsubparagraph}[subparagraph]
\renewcommand{\thesubsubparagraph}{\thesubparagraph.\arabic{subsubparagraph}}
\makeatletter
\providecommand*{\toclevel@subsubparagraph}{6}
\providecommand*{\theHsubsubparagraph}{\thesubparagraph.\arabic{subsubparagraph}}
\providecommand*{\l@subsubparagraph}[2]{}
\makeatother
\titleformat{\subsubparagraph}[block]
{\normalfont\normalsize\bfseries}
{\thesubsubparagraph}{1em}{}

\titlespacing*{\paragraph}{0pt}{\baselineskip}{0.25\baselineskip}
\titlespacing*{\subparagraph}{0pt}{\baselineskip}{0.25\baselineskip}
\titlespacing*{\subsubparagraph}{0pt}{\baselineskip}{0.25\baselineskip}

\newcommand{\titulolargo}{Generación de itinerarios de viaje a partir del procesamiento de correos electrónicos de reservas}
\newcommand{\titulocorto}{TFI - ECD 2024B - Ing. Franco Román Meola}

\usepackage{fancyhdr}
\usepackage{float}
\usepackage{pdflscape}
\usepackage{placeins}

\setlength{\headheight}{14pt}
\pagestyle{fancy}
\fancyhf{}
\fancyhead[C]{\normalfont\titulocorto}
\fancyfoot[C]{\raisebox{-20pt}{\thepage}}
\renewcommand{\headrulewidth}{0pt}
\renewcommand{\footrulewidth}{0pt}
\fancypagestyle{plain}{%
    \fancyhf{}
    \fancyhead[C]{\normalfont\titulocorto}
    \fancyfoot[C]{\raisebox{-20pt}{\thepage}}
    \renewcommand{\headrulewidth}{0pt}
    \renewcommand{\footrulewidth}{0pt}
}
